% =============================================================================
% capitulos/cap1_introducao.tex
%   Cap. 1 -- Introdução e Fundamentação (alvo: 2-3 páginas)
%   Redigido no Passo 3a. Conteúdo de aprofundamento (revisão de literatura
%   estendida) ficou no Apêndice E -- ver pos-textuais/apendices/apendice_e_revisao_estendida.tex.
% =============================================================================
\chapter{Introdução}
\label{cap:introducao}

\section{Contexto}
\label{sec:contexto}

A presença de gatos de vida livre (\textit{Felis catus}) em campi universitários é fenômeno globalmente documentado e exige respostas técnicas que conciliem bem-estar animal, saúde pública e conservação da fauna nativa. No Brasil, o paradigma técnico-jurídico consolidado para o manejo dessas populações é o protocolo \textbf{Captura--Esterilização--Devolução (CED)}, tradução operacional do \textit{Trap--Neuter--Return (TNR)}, cuja eficácia em escala continental e em horizonte de longo prazo é descrita em estudos longitudinais clássicos \cite{levy2003longterm,natoli2019italy,kreisler2019orcat} e atualizada em meta-análises recentes \cite{cecchetti2025planning,schiefer2025tnr}. A literatura é convergente em mostrar que programas TNR de \textbf{alta intensidade espacialmente contíguos} reduzem a densidade populacional e a entrada em abrigos, com ganhos consistentes em bem-estar \cite{gunther2022reduction,spehar2019backtoschool}. Em campus universitário, o estudo de referência é o programa da \textit{University of Central Florida}, que em 28 anos reduziu sua colônia em mais de 85\% mantendo apenas TNR e adoção \cite{spehar2019backtoschool}, padrão arquitetural reproduzido em diversas universidades latino-americanas \cite{junqueira2024evaluation,passos2025population,vecchi2022ethical}.

No \textbf{Campus 2 da USP São Carlos}, a colônia de aproximadamente 50 gatos distribuídos em 10 pontos de alimentação é manejada pela atividade de extensão \textbf{AEX Gatosdoc2}, projeto que reúne docentes, técnicos, discentes e voluntários da comunidade externa. A operação segue os princípios do CED, em conformidade com a \textbf{Lei Federal nº 13.426/2017}~\cite{brasil_lei13426_2017}, a \textbf{Lei Federal nº 15.046/2024}~\cite{brasil_lei15046_2024} e, sobretudo, o \textbf{Decreto Federal nº 12.439/2025}~\cite{brasil_decreto12439_2025}, que institui o Programa Nacional de Proteção e Manejo Populacional Ético de Cães e Gatos e o Cadastro Nacional de Animais Domésticos. A identificação dos animais esterilizados segue a \textbf{Resolução CFMV nº 1.595/2024}~\cite{cfmv_resolucao1595_2024}, que padronizou o corte reto da ponta da orelha (\textit{ear-tipping}) como marcador visual primário, complementada pelas diretrizes operacionais da \textbf{Resolução CFMV nº 1.596/2024}~\cite{cfmv_resolucao1596_2024}.

\section{O problema de engenharia}
\label{sec:problema}

A operação cotidiana do AEX Gatosdoc2 é hoje sustentada por \textbf{vigilância humana presencial}: voluntários realizam visitas regulares aos pontos de alimentação para reposição de ração, observação clínica e tentativa de manter um cadastro atualizado de indivíduos. Esse modelo enfrenta limitações intrínsecas bem documentadas na literatura. Primeiro, a \textbf{cobertura temporal} é necessariamente esparsa, visitas ocorrem em horários compatíveis com a disponibilidade humana, mas a atividade felina é majoritariamente \textbf{crepuscular e noturna}~\cite{cove2022capital,herrera2022prey}. Segundo, a \textbf{identificação individual por inspeção visual} é inviável em populações com pelagens parecidas (sólidos pretos, tabbys cinzas, bicolores brancos) e é o gargalo unanimemente apontado em estudos de \textit{capture--recapture} urbano~\cite{cove2022capital,nakashima2025rest,akbar2025bodypart}. Terceiro, a chegada de \textbf{novos indivíduos abandonados} --- evento crítico no campus, com impacto direto na dispersão da população, depende hoje da percepção subjetiva dos cuidadores, sem registro objetivo do momento ou do ponto de entrada.

O monitoramento sistemático por \textbf{armadilhas fotográficas} (\textit{camera trapping}) é a tradição metodológica madura para esse tipo de questão, com trinta anos de literatura em ecologia de fauna~\cite{bruce2025largescale,cordier2022africa,swanson2015snapshot,cove2022capital}. Câmeras com sensor PIR registram presença sem perturbação humana, permitindo estimar riqueza, ocupação, padrões temporais de atividade e, com identificação individual, densidade populacional. A barreira histórica, porém, é o \textbf{volume de imagens}: projetos de escala média geram dezenas a centenas de milhares de capturas que excedem a capacidade humana de classificação~\cite{norouzzadeh2018automatically,bruce2025largescale}. A virada do \textbf{aprendizado profundo aplicado a camera trapping} resolveu esse gargalo: \citeonline{norouzzadeh2018automatically} demonstraram automação de 99,3\% das classificações do \textit{Snapshot Serengeti} com acurácia equivalente à humana, e o ecossistema atual \emph{MegaDetector + PyTorch-Wildlife}~\cite{hernandez2024pytorchwildlife,microsoft_cameratraps_repo} entrega detecção genérica de animais com licença permissiva e modelos prontos para domínios novos.

A lacuna que motiva este trabalho é específica e atual: \textbf{não há, na literatura publicada, sistema integrado de monitoramento por visão computacional voltado a colônias de gatos comunitários em campus universitário brasileiro}, com o trinômio (i)~detecção robusta a domínios não vistos, (ii)~\textbf{re-identificação individual} (re-ID) capaz de operar em \textit{open-set}, isto é, reconhecer indivíduos novos não cadastrados, e (iii)~validação humana \emph{no laço} para apoio à decisão dos cuidadores. O trabalho metodologicamente mais próximo é o de \citeonline{akbar2025bodypart}, que aplica re-ID por partes do corpo (face, flanco, cauda) a gatos ferais australianos em vegetação natural; o cenário urbano de campus, com diversidade de pelagem e infraestrutura cooperativa de cuidadores, ainda \textbf{não está coberto}.

\section{Objetivos}
\label{sec:objetivos}

\subsection*{Objetivo geral}

\textbf{Estudar e conceber um sistema de monitoramento e apoio ao manejo, baseado em visão computacional, para a colônia de gatos comunitários do Campus 2 da USP São Carlos}, integrando captura local de imagens, processamento por modelos pré-treinados com adaptação ao domínio e identificação individual com validação humana, em conformidade com o protocolo CED e com a legislação brasileira vigente.

\subsection*{Objetivos específicos}

\begin{enumerate}
  \item[\textbf{O1.}] Caracterizar operacionalmente os 10 pontos de alimentação do Campus 2 quanto a energia, conectividade, fluxo de pessoas e fluxo de animais, gerando os requisitos de hardware (Cap.~\ref{cap:hardware}).
  \item[\textbf{O2.}] Comparar arquiteturas candidatas de aquisição de imagem --- IPCAM, ESP32-CAM e Raspberry Pi com NPU --- segundo os critérios de viabilidade, custo total de propriedade e adequação ao campus (Cap.~\ref{cap:hardware}).
  \item[\textbf{O3.}] Conceber o \textbf{pipeline de visão computacional} em quatro fases (ingestão, detecção, classificação, re-identificação), reutilizando modelos abertos de domínio público (MegaDetector v6, CatFLW, PetFace, WildlifeReID-10k) com adaptação local~(Cap.~\ref{cap:pipeline}).
  \item[\textbf{O4.}] Definir o \textbf{protocolo metodológico} de avaliação, com escolha de \textit{datasets} de validação, métricas (mAP, F\textsubscript{1}, Rank-$k$), e desenho experimental para falhas previsíveis (oclusão, baixa luminosidade, indivíduos novos) (Cap.~\ref{cap:metodologia}).
  \item[\textbf{O5.}] \textbf{Documentar a continuidade}, entregando matriz de requisitos funcionais, manuais operacionais e roteiros de instalação que permitam a fase de implementação ser executada por outra equipe de engenharia (Cap.~\ref{cap:conclusao} e Apêndices A--D).
\end{enumerate}

Trata-se, portanto, de um projeto desta natureza de \textbf{estudo e concepção} --- a contribuição central é a \emph{arquitetura} do sistema, a justificativa técnica das escolhas e os \emph{artefatos de engenharia} (especificações, diagramas, protocolos, matriz P01--P10) que viabilizam a implementação posterior. \emph{Implementação plena, instalação e validação em campo permanecem como trabalho futuro.}

Nesse contexto, o extenso detalhamento de infraestrutura física e hardware conduzido no Capítulo 2 não visa à entrega de um protótipo físico finalizado (PoC), mas atua estritamente como ferramenta para derivar os parâmetros operacionais e as restrições reais — como variabilidade geométrica, intermitência energética e presença de fauna não-alvo — que a arquitetura de software necessitará suportar.

\section{Justificativa}
\label{sec:justificativa}

A pertinência do trabalho está em três eixos. \textbf{(i)~Saúde Única}: gatos comunitários estão no centro de questões zoonóticas relevantes no contexto brasileiro --- notadamente a esporotricose felina por \textit{Sporothrix brasiliensis}, hoje em expansão epidêmica no sudeste e sul do país~\cite{ferreira2026tndrtm,bastos2025sporotrichosis,colombo2024emergence}, e a toxoplasmose, com seroprevalência documentada em colônias universitárias brasileiras~\cite{kmetiuk2021toxoplasma,higa2010toxoplasma}. Monitoramento contínuo, sistemático e individualizado é \emph{insumo direto} para vigilância sanitária integrada~\cite{ferroglio2024integrated} e para a tomada de decisão clínica por médicos veterinários~\cite{santiago2025sporothrix}. \textbf{(ii)~Conservação}: predação por gatos sobre fauna nativa é problema documentado globalmente~\cite{loss2013impact,li2021china} e tem revisão específica brasileira recente~\cite{goncalves2025wildcat}; estimar e mapear pressão de predação no campus exige individualização, o que só a re-ID viabiliza em escala. \textbf{(iii)~Lacuna técnica brasileira}: enquanto a Europa dispõe da pilha DeepFaune~\cite{rigoudy2023deepfaune} e os EUA da plataforma Wildlife Insights/SpeciesNet~\cite{wildlifeinsights,google_speciesnet_blog}, \textbf{não existe pilha pública brasileira} de monitoramento por visão computacional para fauna sinantrópica nem, em particular, para gatos comunitários. Este trabalho contribui justamente para fechar essa lacuna, propondo arquitetura aberta e replicável.

A escolha por \textbf{reuso de modelos pré-treinados} ao invés de treinamento do zero é fundamentada pela literatura: detectores generalistas treinados em centenas de datasets, como o MegaDetector, generalizam significativamente melhor para regiões geográficas não vistas do que classificadores multi-espécie treinados localmente~\cite{schneider2018avss,velez2022choosing}. O pipeline proposto (detecção genérica $\rightarrow$ classificação local $\rightarrow$ re-ID individual) replica o padrão de engenharia recomendado por \citeonline{mulero2025addressing}, que atingiu F\textsubscript{1} de 96{,}2\% sobre 1{,}3~milhão de imagens em 24 espécies de mamíferos com poucos milhares de exemplos locais. O módulo de re-ID combina o método de \citeonline{akbar2025bodypart} (\textit{part-based} para gatos ferais) com a verificação face-a-face habilitada pelo dataset \textbf{PetFace}~\cite{shinoda2024petface}, o maior público para identificação individual em pets, com 257~mil indivíduos. A validação cruzada do sistema usará o benchmark \textbf{WildlifeReID-10k}~\cite{adam2024wildlifereid10k} como referência externa, alinhada à competição AnimalCLEF~\cite{animalclef2026}.

\section{Organização do trabalho}
\label{sec:organizacao}

O restante da monografia está organizado em quatro capítulos. O Capítulo~\ref{cap:hardware} trata da \textbf{aquisição de imagem em campo}, discutindo em fluxo corrido a caracterização dos 10 pontos do Campus 2, as três famílias de arquiteturas candidatas, energia e conectividade, e a decisão arquitetural recomendada. O Capítulo~\ref{cap:pipeline} descreve em quatro fases, ingestão, detecção, classificação e re-identificação, como os dados brutos das câmeras são transformados em registros individuais auditáveis, com fundamentação em modelos abertos. O Capítulo~\ref{cap:metodologia} detalha datasets de validação, métricas, desenho experimental e tratamento de falhas previsíveis. O Capítulo~\ref{cap:conclusao} consolida as contribuições, discute condições de contorno e indica trabalhos futuros, com ênfase em continuidade operacional.

Os \textbf{Apêndices} reúnem os artefatos de engenharia produzidos pelo autor: a matriz P01--P10 dos pontos de alimentação (Apêndice~\ref{ap:matriz-pontos}) e os códigos do pipeline (Apêndice~\ref{ap:pseudocodigos}). Os demais artefatos operacionais --- requisitos funcionais, manuais e \emph{datasheets} de referência --- são mantidos como documentação de execução no repositório do projeto, fora desta monografia.